% funding-rate.tex — Perpetual funding rate (architecture-independent properties)
% Kontrol: prove_cfmm_perpetual_fundingConvergence

\section{Funding Rate}\label{sec:funding-rate}

This instrument is perpetual: it has no fixed expiry but settles at liquidity events.
The funding rate mechanism replaces settlement, following the framework of
\cite{angeris2022perpetuals}.

\subsection{Mark and Index Prices}

\begin{definition}[Index Price]\label{def:index-price}
The index price is derived from the hook's co-primary state:
\begin{equation}\label{eq:index-price}
  p_{\text{index}} = \frac{\DeltaPlus}{1 - \DeltaPlus}
\end{equation}
where $\DeltaPlus = \max(0,\, \AT - \ATnull)$ is the realized concentration deviation
from the underlying pool. This is computed on the fly from stored $\AT$, $\ThetaSum$, $\PosCount$.
\end{definition}

\begin{definition}[Mark Price]\label{def:mark-price}
The mark price is the protection mechanism's implied price of fee concentration.
Its definition depends on the architecture choice (\cref{sec:architecture-pending}):
\begin{itemize}
  \item \textbf{If CFMM:} $p_{\text{mark}}$ is the spot price implied by the trading function.
  \item \textbf{If hook insurance:} $p_{\text{mark}}$ is the implied cost of protection
    derived from the fee extraction rate $\gamma$ and accumulated reserve.
\end{itemize}
\end{definition}

\subsection{Discrete Funding at Liquidity Events}

Funding is exchanged at each liquidity event (\texttt{afterAddLiquidity} or
\texttt{afterRemoveLiquidity} in the underlying pool), not continuously.

\begin{definition}[Funding Payment]\label{def:funding-payment}
At liquidity event $n$, the funding payment per unit of exposure is:
\begin{equation}\label{eq:funding-payment}
  F_n = \fundrate_n \cdot (p_{\text{mark},n} - p_{\text{index},n})
\end{equation}
where $\fundrate_n > 0$ is the funding rate coefficient at event $n$.
\end{definition}

\subsection{Jump-Adjusted Replication}

\begin{theorem}[Jump Safety]\label{thm:jump-safety}
When $\DeltaPlus$ jumps discretely from $\delta$ to $\delta'$ at a liquidity event
(due to abrupt changes in $\PosCount$ or $\ThetaSum$), the protection value change
is bounded:
\begin{equation}\label{eq:jump-bound}
  |\Delta\text{value}| \leq |V(\delta') - V(\delta)|
\end{equation}
where $V$ is the (monotone) protection payoff function.
\end{theorem}

\begin{proof}[Proof sketch]
$\DeltaPlus$ can jump when a large position mints (increasing $\PosCount$, potentially
decreasing $\ATnull$ and thus increasing $\DeltaPlus$) or burns (decreasing $\PosCount$).
The discrete funding mechanism compensates for the discontinuity, following
\cite[Section~3]{angeris2022perpetuals}. The bound holds because the oracle-free
construction adjusts to the new state after each event, and the payoff's
monotonicity ensures the value change is bounded by the payoff difference. \qedhere
\end{proof}

\subsection{Convergence Property}

\begin{proposition}[Mark-Index Convergence]\label{prop:convergence}
Under the funding mechanism, $|p_{\text{mark}} - p_{\text{index}}|$ decreases
after each liquidity event (in expectation).
\end{proposition}

\begin{proof}[Proof sketch]
When $p_{\text{mark}} > p_{\text{index}}$, the funding cost discourages further
mark appreciation. The converse holds when $p_{\text{mark}} < p_{\text{index}}$.
A formal convergence proof under specific assumptions is deferred to the
Kontrol verification phase. \qedhere
\end{proof}
